\documentclass[11pt, letterpaper]{article}

% --- Essential Packages ---
\usepackage[utf8]{inputenc}
\usepackage[T1]{fontenc}
\usepackage{lmodern}
\usepackage{mathpazo} % Palatino font for a professional look
\usepackage{amsmath, amssymb, amsthm}
\usepackage{geometry}
\usepackage{booktabs}
\usepackage{graphicx}
\usepackage{setspace}
\usepackage[dvipsnames]{xcolor}
\usepackage{hyperref}

% --- Page Setup ---
\geometry{margin=1.25in}
\setstretch{1.15}
\hypersetup{
    colorlinks=true,
    linkcolor=MidnightBlue,
    citecolor=MidnightBlue,
    urlcolor=MidnightBlue
}

% --- Theorem Environments ---
\newtheorem{theorem}{Theorem}
\newtheorem{lemma}{Lemma}
\newtheorem{proposition}{Proposition}
\newtheorem{assumption}{Assumption}
\newtheorem{corollary}{Corollary}
\theoremstyle{remark}
\newtheorem{remark}{Remark}

% --- Custom Math Operators ---
\DeclareMathOperator{\E}{\mathbb{E}}
\DeclareMathOperator{\Var}{Var}
\DeclareMathOperator{\Cov}{Cov}
\DeclareMathOperator{\plim}{plim}
\newcommand{\Omegahat}{\widehat{\Omega}}
\newcommand{\lambdahat}{\widehat{\lambda}}

% --- Title Information ---
\title{A Continuous-Dose Extension of the Causal Cluster Variance Estimator: Theory and Proofs}
\author{Formal Econometric Derivation}
\date{\today}

\begin{document}

\maketitle

\begin{abstract}
    \noindent We extend the fixed-effect Causal Cluster Variance (CCV) estimator proposed by Abadie, Athey, Imbens, and Wooldridge (2023) from a binary treatment framework to continuous-dose treatments. By defining a cluster-level continuous variation parameter based on Frisch-Waugh-Lovell residualized treatment vectors, we derive a continuous-treatment CCV weight that forms a strictly convex combination of the Eicker-Huber-White robust variance and the Liang-Zeger clustered variance estimators. We formally prove that this formula is a strict mathematical generalization of the AAIW binary weight. Furthermore, we define the higher-moment location-scale assumptions required for asymptotic optimality and demonstrate the estimator's boundary behavior using the Adverse Effect Wage Rate (AEWR) ``bite'' as a censored continuous application.
\end{abstract}

\vspace{2em}
\hrule
\vspace{2em}

\section{Setup and The Residual-Maker Projection}

Let $g=1,\dots,G$ index counties or county-like clusters, and let $i=1,\dots,n_g$ index observations within county $g$ (for example, county-year, worker-county-year, or farm-county-year observations). Let $Y_{gi}$ be the observed outcome, and let $D_{gi}$ be the continuous treatment intensity. An applied example of $D_{gi}$ is the continuous Adverse Effect Wage Rate (AEWR) ``bite'' intensity:
\begin{equation}
    D_{gi} = \max\{0, \log(\text{AEWR}_{gi}) - \log(w^0_{gi})\}.
\end{equation}
Let $Z_{gi}$ collect all covariates and fixed effects (e.g., county fixed effects, year fixed effects, control variables) other than $D_{gi}$. The regression model of interest is:
\begin{equation}
    Y_{gi} = \beta D_{gi} + Z_{gi}'\theta + u_{gi}.
\end{equation}
Stacking the data for the full sample of size $N = \sum_{g=1}^G n_g$, we write the system as $Y = D\beta + Z\theta + u$.

Let $M_Z = I - Z(Z'Z)^{-1}Z'$ be the residual-maker matrix after projecting onto the covariates and fixed effects. By the Frisch-Waugh-Lovell (FWL) theorem, we define the residualized variables as:
\begin{equation}
    \widetilde Y = M_Z Y, \qquad \widetilde D = M_Z D, \qquad \widetilde u = M_Z u.
\end{equation}
The Ordinary Least Squares (OLS) estimator for the treatment parameter is given by:
\begin{equation}
    \widehat\beta = \frac{\widetilde D'\widetilde Y}{\widetilde D'\widetilde D}.
\end{equation}
Let $S = \sum_{g=1}^G \sum_{i=1}^{n_g} \widetilde D_{gi}^2$. Substituting $\widetilde Y = \widetilde D\beta + \widetilde u$, the estimation error expands to:
\begin{equation}
    \widehat\beta - \beta = S^{-1} \sum_{g=1}^G \sum_{i=1}^{n_g} \widetilde D_{gi}\widetilde u_{gi}.
\end{equation}
Defining the observation-level score as $\psi_{gi} = \widetilde D_{gi}\widetilde u_{gi}$ and the county-level score as $\Psi_g = \sum_{i=1}^{n_g}\psi_{gi}$, we rewrite the estimation error as:
\begin{equation}
    \widehat\beta - \beta = S^{-1}\sum_{g=1}^G \Psi_g. \label{eq:beta_error}
\end{equation}
The variance of this object is the primary target of our estimation strategy.

\section{Robust vs. Clustered Variance Estimators}

The difference between standard error regimes is rooted in the treatment of the cross-products within equation \eqref{eq:beta_error}. The Eicker-Huber-White (robust) variance estimator treats the observation-level scores independently:
\begin{equation}
    \widehat V_{\text{robust}} = S^{-2} \sum_{g=1}^G \sum_{i=1}^{n_g} \widehat\psi_{gi}^2.
\end{equation}
Conversely, the standard Liang-Zeger cluster-robust variance estimator aggregates scores at the cluster level before squaring:
\begin{equation}
    \widehat V_{\text{cluster}} = S^{-2} \sum_{g=1}^G \left( \sum_{i=1}^{n_g}\widehat\psi_{gi} \right)^2.
\end{equation}

\begin{lemma}
The difference between the cluster-robust and heteroskedasticity-robust variance estimators is solely comprised of within-cluster score cross-products.
\end{lemma}
\begin{proof}
Expanding the squared sum in $\widehat V_{\text{cluster}}$ yields:
\begin{align*}
    \widehat V_{\text{cluster}} &= S^{-2} \sum_{g=1}^G \left[ \sum_{i=1}^{n_g}\widehat\psi_{gi}^2 + \sum_{i\neq j} \widehat\psi_{gi}\widehat\psi_{gj} \right] \\
    &= \widehat V_{\text{robust}} + S^{-2} \sum_{g=1}^G \sum_{i\neq j} \widehat\psi_{gi}\widehat\psi_{gj}.
\end{align*}
\end{proof}

According to the design-based framework established by Abadie, Athey, Imbens, and Wooldridge (2023) \cite{AAIW2023}, the inclusion of the cross-product term $\sum_{i \neq j} \widehat\psi_{gi}\widehat\psi_{gj}$ should be governed by the fraction of the treatment assignment variation that occurs across clusters versus within clusters.

\section{The Continuous-Treatment Causal Cluster Variance Estimator}

In the AAIW (2023) framework with a binary treatment $W_{gi} \in \{0,1\}$ and cluster fixed effects, the treatment variation parameter is based strictly on the within-cluster variance of the binary indicator, $\bar W_g(1-\bar W_g)$.

To generalize this logic to a continuous treatment vector $D$, we define the continuous analog using the FWL-residualized treatment vector. Let the residualized treatment variance within cluster $g$ be:
\begin{equation}
    \widehat\Omega_g = \frac{1}{n_g} \sum_{i=1}^{n_g} \widetilde D_{gi}^2.
\end{equation}
If panel lengths $n_g$ are highly unequal and the target estimand is unit-weighted rather than equal-county-weighted, we define the sum-of-squares metric:
\begin{equation}
    \widehat S_g = \sum_{i=1}^{n_g} \widetilde D_{gi}^2.
\end{equation}

\begin{theorem}[Continuous CCV Estimator]
Let $q \in [0,1]$ represent the population sampling fraction. The continuous-dose causal cluster variance weight is defined as:
\begin{equation}
    \widehat\lambda_D = 1 - q \frac{\left(G^{-1}\sum_{g=1}^G \widehat\Omega_g\right)^2}{G^{-1}\sum_{g=1}^G \widehat\Omega_g^2}.
\end{equation}
The final variance estimator is the resulting convex combination:
\begin{equation}
    \widehat V_{\text{CCV},D} = \widehat\lambda_D \widehat V_{\text{cluster}} + (1-\widehat\lambda_D)\widehat V_{\text{robust}}.
\end{equation}
\end{theorem}
For unbalanced panel specifications targeting the unit-level Average Treatment Effect, the estimator substitutes $\widehat S_g$ for $\widehat\Omega_g$:
\begin{equation}
    \widehat\lambda_D^{\text{unit}} = 1 - q \frac{\left(G^{-1}\sum_{g=1}^G \widehat S_g\right)^2}{G^{-1}\sum_{g=1}^G \widehat S_g^2}.
\end{equation}

\section{Proof of Reduction to the AAIW Binary Case}

We now prove that the continuous weight strictly maps to the AAIW formula in the boolean domain.

\begin{proposition}
If $D_{gi} = W_{gi} \in \{0,1\}$ and covariates $Z_{gi}$ consist strictly of cluster fixed effects, the continuous variation term $\widehat\Omega_g$ reduces exactly to the AAIW fixed-effect parameter $\bar W_g(1-\bar W_g)$.
\end{proposition}

\begin{proof}
Because $Z_{gi}$ contains only cluster fixed effects, the FWL residualized treatment simplifies to the within-cluster demeaned treatment:
\begin{equation*}
    \widetilde D_{gi} = W_{gi} - \bar W_g.
\end{equation*}
Therefore, our continuous parameter $\widehat\Omega_g$ is:
\begin{equation*}
    \widehat\Omega_g = \frac{1}{n_g} \sum_{i=1}^{n_g} (W_{gi} - \bar W_g)^2.
\end{equation*}
Expanding the squared binomial yields:
\begin{equation*}
    \widehat\Omega_g = \frac{1}{n_g} \sum_i \left(W_{gi}^2 - 2W_{gi}\bar W_g + \bar W_g^2\right).
\end{equation*}
Because $W_{gi}$ is strictly binary, we rely on the fundamental identity $W_{gi}^2 = W_{gi}$. Substituting this identity gives:
\begin{align*}
    \widehat\Omega_g &= \frac{1}{n_g} \sum_i W_{gi} - 2\bar W_g \left( \frac{1}{n_g}\sum_i W_{gi} \right) + \bar W_g^2 \\
    &= \bar W_g - 2\bar W_g^2 + \bar W_g^2 \\
    &= \bar W_g(1-\bar W_g).
\end{align*}
Plugging this directly into the continuous formula establishes the exact equivalence:
\begin{equation*}
    \widehat\lambda_D = 1 - q \frac{\left(G^{-1}\sum_g \bar W_g(1-\bar W_g)\right)^2}{G^{-1}\sum_g \bar W_g^2(1-\bar W_g)^2},
\end{equation*}
which identically matches the AAIW (2023) fixed-effect CCV weight. Thus, the continuous formula operates as a strict generalization of the original framework.
\end{proof}

\section{Proof of Valid Convexity}

To ensure $\widehat V_{\text{CCV},D}$ behaves dynamically without breaking finite sample bounds, we prove $\widehat\lambda_D$ represents a valid probability measure for interpolation.

\begin{theorem}
For any finite continuous treatment distribution where $q \in [0,1]$ and not all $\widehat\Omega_g = 0$, $0 \leq \widehat\lambda_D \leq 1$.
\end{theorem}
\begin{proof}
Let $\Omegahat \in \mathbb{R}^G$ be a vector of strictly non-negative real numbers representing the cluster variation parameters. By the Cauchy-Schwarz inequality, the inner product of the vector $\Omegahat$ with a vector of ones $\mathbf{1}$ satisfies:
\begin{equation*}
    \langle \Omegahat, \mathbf{1} \rangle^2 \leq \langle \Omegahat, \Omegahat \rangle \langle \mathbf{1}, \mathbf{1} \rangle.
\end{equation*}
Expanding the sums, this requires:
\begin{equation*}
    \left( \sum_{g=1}^G \widehat\Omega_g \right)^2 \leq G \sum_{g=1}^G \widehat\Omega_g^2.
\end{equation*}
Dividing both sides by $G^2$ yields the expectation limits:
\begin{equation*}
    \left( G^{-1}\sum_{g=1}^G \widehat\Omega_g \right)^2 \leq G^{-1}\sum_{g=1}^G \widehat\Omega_g^2.
\end{equation*}
Because both the numerator and denominator are strictly non-negative, taking their ratio enforces:
\begin{equation*}
    0 \leq \frac{ \left( G^{-1}\sum_g \widehat\Omega_g \right)^2 }{ G^{-1}\sum_g \widehat\Omega_g^2 } \leq 1.
\end{equation*}
Because $q \in [0,1]$, multiplying the fraction by $q$ preserves the bounds $0 \leq q(\dots) \leq 1$. Thus, subtracting the term from $1$ guarantees:
\begin{equation*}
    0 \leq 1 - q \frac{\left(G^{-1}\sum_g \widehat\Omega_g\right)^2}{G^{-1}\sum_g \widehat\Omega_g^2} \leq 1.
\end{equation*}
Therefore, $\widehat\lambda_D \in [0,1]$, confirming $\widehat V_{\text{CCV},D}$ is a strict convex combination.
\end{proof}

\section{Consistency Framework}

We formalize the asymptotic limits using standard regularity assumptions under Slutsky's theorem.

\begin{assumption}
As the number of clusters grows ($G \to \infty$), the relative size of any single cluster's contribution remains bounded, such that $\max_g \frac{S_g}{S} \to 0$. Furthermore, the residualized treatment and errors exhibit bounded fourth moments, and the empirical variance estimators possess non-zero limiting behavior:
\begin{align*}
    G^{-1}\sum_g \widehat\Omega_g &\xrightarrow{p} \Omega_1 > 0, \\
    G^{-1}\sum_g \widehat\Omega_g^2 &\xrightarrow{p} \Omega_2 > 0.
\end{align*}
\end{assumption}

Under Assumption 1, the Continuous Probability Mapping Theorem and Slutsky's Theorem immediately give:
\begin{equation*}
    \widehat\lambda_D \xrightarrow{p} \lambda_D = 1 - q \frac{\Omega_1^2}{\Omega_2}.
\end{equation*}
Given that under the standard cluster-level law of large numbers, $\widehat V_{\text{robust}} \xrightarrow{p} V_{\text{robust}}$ and $\widehat V_{\text{cluster}} \xrightarrow{p} V_{\text{cluster}}$, we conclude:
\begin{equation*}
    \widehat V_{\text{CCV},D} \xrightarrow{p} \lambda_D V_{\text{cluster}} + (1-\lambda_D)V_{\text{robust}}.
\end{equation*}

\section{The Distributional Limit: Higher-Order Moments}

The critical caveat in continuous generalization stems from higher-order statistical moments. As established in Proposition 1, a binary treatment enjoys the algebraic property $W^2 = W$, which systematically collapses higher moments. Thus, terms like $E[W^4]$ naturally compress into functions of the mean, allowing the second-moment parameter $\bar{W}_g(1-\bar{W}_g)$ to fully capture the shape of the design variance.

For an arbitrary continuous treatment $D$, the optimality of the design variance relies heavily on $E[\widetilde D^3 \mid g]$ and $E[\widetilde D^4 \mid g]$, which are no longer intrinsically anchored to $\Omega_g$.

\begin{theorem}[Location-Scale Sufficiency]
The parameter $\Omega_g$ serves as an asymptotically sufficient statistic for the optimal continuous design weight if the treatment assignment process satisfies a common-shape location-scale restriction.
\end{theorem}
\begin{proof}
Assume the residualized continuous treatment is generated independently given $g$ by:
\begin{equation*}
    \widetilde D_{gi} = \sigma_g \epsilon_{gi},
\end{equation*}
where $\epsilon_{gi}$ is an identically distributed random error with mean 0, variance 1, and global kurtosis $\kappa < \infty$. Consequently:
\begin{equation*}
    \widehat\Omega_g = \frac{1}{n_g}\sum_{i} \widetilde D_{gi}^2 \xrightarrow{p} \sigma_g^2.
\end{equation*}
The fourth moment evaluates as:
\begin{equation*}
    E[\widetilde D_{gi}^4 \mid g] = \sigma_g^4 E[\epsilon_{gi}^4] = (\sigma_g^2)^2 \kappa \equiv \Omega_g^2 \kappa.
\end{equation*}
Because the higher moments scale proportionally with $\Omega_g^2$, $\Omega_g$ sufficiently summarizes the cluster-level identifying variation in the assignment process. If kurtosis is heavily heterogeneous across clusters, the CCV estimator remains a computationally valid convex combination but may deviate from strict asymptotic efficiency.
\end{proof}

\section{Application: The AEWR Zero-Bite Boundary}

Applying this estimator to the Adverse Effect Wage Rate (AEWR) highlights a desirable empirical property. The treatment variable is left-censored at zero:
\begin{equation*}
    D_{ct} = \max\{0, \log(\text{AEWR}_{rt}) - \log(w^0_{ct})\}.
\end{equation*}
In high-wage agricultural counties, $w^0_{ct}$ chronically exceeds the statutory $\text{AEWR}_{rt}$, resulting in $D_{ct} = 0$ across all time periods $t$. Consequently, after residualization, $\widetilde{D}_{ct} = 0$, yielding $\widehat\Omega_c = 0$.

If treatment is clustered strictly in a handful of low-wage counties, the proportion of zeros dominates the mean absolute variance. As the sum of squares $G^{-1}\sum \widehat\Omega_c^2$ remains anchored by the few highly treated clusters, the ratio strictly vanishes:
\begin{equation*}
    \frac{\left(G^{-1}\sum \widehat\Omega_c\right)^2}{G^{-1}\sum \widehat\Omega_c^2} \longrightarrow 0 \quad \implies \quad \widehat\lambda_D \longrightarrow 1.
\end{equation*}
Therefore, $\widehat V_{\text{CCV},D} \to \widehat V_{\text{cluster}}$. This boundary limit is highly attractive: if a policy variation is heavily concentrated in a few discrete locations, design-based inference insists that treatment assignment behaves essentially at the macro-level, requiring fully clustered standard errors. The $\widehat\Omega_g$ extension identifies this boundary and collapses optimally without manual intervention.

\vspace{2em}
\begin{thebibliography}{9}
\bibitem{AAIW2023}
Abadie, A., Athey, S., Imbens, G. W., \& Wooldridge, J. M. (2023).
When Should You Adjust Standard Errors for Clustering?
\textit{The Quarterly Journal of Economics}, 138(1), 1--35.
\end{thebibliography}

\end{document}
